\documentclass[11pt]{article}
\usepackage[a4paper,margin=24mm]{geometry}
\usepackage{amsmath,amssymb,amsthm,mathtools,bm,slashed}
\usepackage{booktabs,longtable,array}
\usepackage{enumitem}
\usepackage{xcolor}
\usepackage{microtype}
\usepackage[hidelinks]{hyperref}

\setlength{\emergencystretch}{3em}
\newtheorem{theorem}{Theorem}[section]
\newtheorem{proposition}[theorem]{Proposition}
\newtheorem{lemma}[theorem]{Lemma}
\newtheorem{corollary}[theorem]{Corollary}
\theoremstyle{definition}
\newtheorem{definition}[theorem]{Definition}
\newtheorem{gate}[theorem]{Fail-closed gate}
\theoremstyle{remark}
\newtheorem{remark}[theorem]{Remark}

\newcommand{\ii}{\mathrm i}
\newcommand{\dd}{\mathrm d}
\newcommand{\CC}{\mathbb C}
\newcommand{\ZZ}{\mathbb Z}
\newcommand{\CP}{\mathbb {CP}}
\newcommand{\cO}{\mathcal O}
\newcommand{\cH}{\mathcal H}
\newcommand{\cM}{\mathcal M}
\newcommand{\cC}{\mathcal C}
\newcommand{\cG}{\mathcal G}
\newcommand{\Tr}{\operatorname{Tr}}
\newcommand{\Ker}{\operatorname{Ker}}
\newcommand{\ind}{\operatorname{ind}}
\newcommand{\status}[1]{\textcolor{blue!60!black}{\textsf{#1}}}
\newcolumntype{L}[1]{>{\raggedright\arraybackslash}p{#1}}

\title{AP-E4: A Physical Tangent Fermion from Vortex Moduli-Space SQM\\
Canonical Spin\texorpdfstring{\({}^c\)}{c} Quantization, Conditional
\(\cO(2)\) Twisting,\\ Chirality, Partner Spectrum, Gap, and the AP-E3
Composability Gate}
\author{Route F research note}
\date{16 July 2026}

\begin{document}
\maketitle

\begin{abstract}
This note supplies the missing intrinsic tangent-fermion realization behind
AP-E4, while showing that the vacuum-line polarization remains independent
physical data.
A half-BPS non-Abelian \(U(2)\) Chern--Simons--Higgs vortex has a normalizable
internal moduli space
\(SU(2)_{c+f}/U(1)\simeq\CP^1\).  Promoting its orientation to a slow
worldline coordinate yields the exact Fubini--Study kinetic term
\[
 L_B=C\frac{|\dot w|^2}{(1+|w|^2)^2},\qquad
 C=\frac{2\pi}{e^2}I_0>0,
\]
where \(I_0\) is an explicitly convergent radial \(L^2\) integral.  The one
complex unbroken supercharge supplies a complex fermion satisfying
\(\psi^v=-w^{-2}\psi^w\) under \(v=1/w\), so the physical worldline fermion is
intrinsically tangent-valued.  This closes the ``bosonic fluctuation is not a
fermion'' objection for this mother model.

The decisive distinction is then made explicit.  In canonical
Spin\({}^c\)/Dolbeault polarization the tangent fermion generates the
Clifford module \(\Omega^{0,*}(\CP^1)\); it does \emph{not} also make the
wavefunction \(T\CP^1=\cO(2)\)-valued.  Consequently the declared canonical
Spin\({}^c\) re-quantization of the untwisted complex has one Dolbeault ground
state, not three; the source-selected half-form/operator ordering instead has
no normalizable zero mode.  An independent coefficient line
\(E=\cO(2)\) changes the supercharge to
\(Q_E=\sqrt2\,\bar\partial_E\) and gives
\((h^0,h^1)=(3,0)\).  Its exact nonzero tower is
\[
 \lambda_{n,\pm}=\pm\frac{2}{\sqrt C}\sqrt{n(n+3)},\qquad
 \operatorname{mult}(\lambda_{n,\pm})=2n+3,\qquad n\ge1,
\]
so \(\Delta_D=4/\sqrt C\) and
\(\Delta_H=8/C\) for \(H=D^2/2\).

AP-E3's oriented level \(k=+2\) mixed-WZW term could provide precisely this
coefficient line, but only after proving that its charged-two-colour
\(B=1\) soliton has the \emph{same} \(\CP^1\) collective-coordinate map and
that the spatially transgressed WZW line pulls back with Chern number \(+2\).
The vortex used here is
an independent \(2+1\)-dimensional mother model, not that soliton.  Thus
\texttt{physical\_tangent\_fermion\_in\_intrinsic\_sqm\_derived=true}, while
the charged-two-colour embedding, AP-E4 worldline completion, Route-E portal,
and physics promotion remain false.  Product compactification is retained
only as a backup behind a full six-dimensional anomaly gate.
\end{abstract}

\tableofcontents

\section{Problem and theorem-level answer}

The AP-E4 question contains four logically distinct layers:
\begin{enumerate}[label=(E4.\arabic*)]
\item derive a physical \(\CP^1\) collective-coordinate manifold and its
  \(L^2\) metric from a mother action;
\item derive a Grassmann-odd tangent partner and the worldline supercharge;
\item determine which additional line bundle, if any, twists the quantum
  states and then solve chirality, partners, and the gap;
\item prove that the result composes with AP-E3 and eventually Route E.
\end{enumerate}
The first two layers are solved by the intrinsic vortex model below.  The
third is solved mathematically both without and with a declared \(\cO(2)\)
coefficient line.  The fourth remains a sharply stated physical gate.

\begin{theorem}[Intrinsic tangent result and fixed-polarization line-twist result]
For a charge-one half-BPS \(U(2)\) non-Abelian Chern--Simons vortex with the
translational coordinate factored out, the internal moduli are \(\CP^1\),
their \(L^2\) metric is Fubini--Study with a finite positive coefficient
\(C\), and the BPS collective-coordinate supersymmetry supplies a physical
tangent fermion.  In the declared canonical Dolbeault re-quantization the untwisted
ground-state space is
\(H^{0,*}(\CP^1,\cO)=\CC\).  Within \(SU(2)\)-equivariant holomorphic line
twists of this fixed canonical operator, and with no added superpotential or
extra Fermi multiplet, exactly three positive-chirality ground states occur
if and only if the independent coefficient line is \(E=\cO(2)\):
\[
 H^{0,*}(\CP^1,E)=H^0(\CP^1,\cO(2))\simeq\CC^3,
 \qquad H^1(\CP^1,\cO(2))=0.
\]
Every nonzero level is paired.  This theorem does not claim uniqueness among
all possible SQM field contents, potentials, or vacuum-line regularizations.
\end{theorem}

\section{A BPS mother model with internal \texorpdfstring{\(\CP^1\)}{CP1}}

\subsection{Bulk ingredients and colour--flavour locking}

Consider the \(2+1\)-dimensional \(\mathcal N=2\) supersymmetric
Chern--Simons--Higgs theory with gauge group \(U(2)\), two fundamental Higgs
flavours assembled into a \(2\times2\) matrix \(H\), and a positive FI scale
\(\xi\).  Its bosonic action has the schematic but sufficient form
\begin{equation}
 S_{\rm bulk}=\int\dd^3x\left[
 \frac{\kappa}{4\pi}\epsilon^{\mu\nu\rho}
 \Tr\!\left(A_\mu\partial_\nu A_\rho-
 \frac{2\ii e}{3}A_\mu A_\nu A_\rho\right)
 +\Tr(D_\mu H)^\dagger D^\mu H-V_6(H;\xi,\kappa,e)
 +\mathcal L_F\right].
 \label{eq:bulk-action}
\end{equation}
Supersymmetry fixes the sextic potential \(V_6\) and admits first-order BPS
vortex equations.  We use the explicit BPS family and collective-coordinate
reduction of Ref.~\cite{AldrovandiSchaposnik2007APE4SQM}; no claim below
depends on inventing a new radial solution.

The asymmetric vacuum may be chosen as
\begin{equation}
 H_\infty=\sqrt\xi\,\bm 1_2.
 \label{eq:locked-vacuum}
\end{equation}
Gauge \(SU(2)_c\) and flavour \(SU(2)_f\) are then locked to a global
diagonal \(SU(2)_{c+f}\).  A reference unit vortex embeds its winding in one
colour--flavour direction and has stabilizer \(U(1)\).  Therefore its
internal orbit is
\begin{equation}
 \cM_{\rm int}=\frac{SU(2)_{c+f}}{U(1)}\simeq\CP^1.
 \label{eq:orbit}
\end{equation}
The full one-vortex moduli space contains also its centre
\(Z\in\CC\): \(\cM_1\simeq\CC_Z\times\CP^1\).  We factor out \(Z\), or
equivalently hold the vortex centre fixed with a weak supersymmetric trap.
This avoids confusing translational states with the proposed three-state
internal sector.

\subsection{Projector representation and physical zero mode}

Write a normalized orientation and its rank-one projector as
\begin{equation}
 n\in\CC^2,\qquad n^\dagger n=1,\qquad
 n\sim e^{\ii\alpha}n,\qquad P=nn^\dagger.
 \label{eq:n-projector}
\end{equation}
The oriented vortex fields can be expressed schematically as
\begin{equation}
 H(r;n)=\sqrt\xi\,\bm1_2+
 \bigl(\widetilde\varphi(r)-\sqrt\xi\bigr)P,
 \qquad
 A_\mu(r;n)=A_\mu^{(0)}(r)\bm1_2+A_\mu^{(1)}(r)
 \left(P-\frac12\bm1_2\right).
 \label{eq:oriented-fields}
\end{equation}
An infinitesimal global rotation generated by a Hermitian
\(\epsilon\in\mathfrak{su}(2)\) gives
\begin{equation}
 \delta_\epsilon P=\ii[\epsilon,P],\qquad
 \delta_\epsilon H=\ii[\epsilon,H],\qquad
 \delta_\epsilon A_\mu=\ii[\epsilon,A_\mu].
 \label{eq:symmetry-zero-mode}
\end{equation}
Differentiating a family of exact solutions along an exact symmetry orbit
annihilates the linearized field equations.  Directions satisfying
\([\epsilon,P]=0\) are the stabilizer and hence zero; the remaining two real
directions span \(T_P\CP^1\).

Normalizability follows from the local-vortex boundary conditions.  Near
the core all profiles are regular, so the radial norm has a bounded
integrand times \(r\dd r\).  At infinity the massive bulk equations imply
\(\widetilde\varphi-\sqrt\xi=O(e^{-mr})\) and analogous decay for the
non-Abelian gauge profile.  Hence
\begin{equation}
 \|\delta_\epsilon\Phi\|_{L^2}^2
 \leq C_0\int_0^1r\,\dd r+
 C_\infty\int_1^\infty r e^{-2mr}\,\dd r<\infty.
 \label{eq:normalizable}
\end{equation}
This is a genuine normalizable collective mode, not merely a tangent vector
in field-value space.  Index and moduli-matrix analyses of non-Abelian
vortices give the same moduli count in broader Yang--Mills--Chern--Simons
settings~\cite{GudnasonEto2021APE4SQM}.

\section{First-principles \texorpdfstring{\(L^2\)}{L2} metric}

\subsection{Horizontal geometry of the normalized doublet}

For a path \(n(t)\), the physical horizontal velocity is
\begin{equation}
 \eta=(1-P)\dot n,\qquad n^\dagger\eta=0,\qquad
 \|\eta\|^2=\dot n^\dagger\dot n-|n^\dagger\dot n|^2.
 \label{eq:horizontal-velocity}
\end{equation}
Because \(n^\dagger n=1\), \(n^\dagger\dot n\) is purely imaginary and
therefore
\begin{equation}
 \dot n^\dagger\dot n+(n^\dagger\dot n)^2
 =\dot n^\dagger\dot n-|n^\dagger\dot n|^2=\|\eta\|^2.
 \label{eq:source-bracket}
\end{equation}
On the north chart choose
\begin{equation}
 n_N(w)=\frac{(1,w)^T}{\sqrt{1+|w|^2}}.
 \label{eq:north-section}
\end{equation}
Direct differentiation and projection give
\begin{equation}
 \|(1-P)\dot n_N\|^2
 =\frac{|\dot w|^2}{(1+|w|^2)^2}.
 \label{eq:fs-shape}
\end{equation}
This proves the Fubini--Study shape without invoking symmetry uniqueness.

\subsection{Radial norm and elimination of the temporal compensator}

Promoting \(n\) to \(n(t)\) requires a temporal gauge-field compensator
with a radial function \(\rho(r)\).  Integrating the explicit bulk ansatz
over the transverse plane gives~\cite{AldrovandiSchaposnik2007APE4SQM}
\begin{align}
 S_{\rm eff}
 &=\frac{2\pi}{e^2}I_0\int\dd t\,
 \left[\dot n^\dagger\dot n+(n^\dagger\dot n)^2\right],
 \label{eq:published-effective-action}\\
 I_0&=\int_0^\infty\dd\hat r\,\hat r\left[
 (\hat\varphi^2+\hat\xi)\rho^2
 +2(1-\rho)(\hat\varphi-\sqrt{\hat\xi})^2\right],
 \label{eq:i0}\\
 \rho&=\frac{(\hat\varphi-\sqrt{\hat\xi})^2}
 {\hat\varphi^2+\hat\xi}.
 \label{eq:rho}
\end{align}
Here hats denote the dimensionless radial variables used in that reduction.
The last equation is algebraic because \(\rho\) is nondynamical in the pure
Chern--Simons model.  Substitution makes finiteness manifest from the same
core regularity and exponential tail used in Eq.~\eqref{eq:normalizable};
the integrand is nonnegative and nonzero for a nontrivial vortex, so
\begin{equation}
 C:=\frac{2\pi}{e^2}I_0>0.
 \label{eq:C}
\end{equation}
Combining Eqs.~\eqref{eq:source-bracket}--\eqref{eq:C} yields
\begin{equation}
 \boxed{L_B=C\frac{\dot w\dot{\bar w}}{(1+|w|^2)^2}},\qquad
 K=C\log(1+|w|^2),\qquad
 g_{w\bar w}=\frac{C}{(1+|w|^2)^2}.
 \label{eq:LB}
\end{equation}
Since a round sphere of radius \(R\) has
\(\dd s^2=4R^2\dd w\dd\bar w/(1+|w|^2)^2\),
\begin{equation}
 R=\frac{\sqrt C}{2},\qquad \operatorname{Area}(\CP^1)=\pi C.
 \label{eq:radius-area}
\end{equation}

For deterministic regression, the companion code also inserts the analytic
surrogate \(\hat\varphi=\tanh\hat r\), \(\hat\xi=1\) into
Eqs.~\eqref{eq:i0}--\eqref{eq:rho}, and verifies positivity and quadrature
convergence.  This is a numerical check of the reduction, not a claim to
have solved the microscopic BPS radial equation.

\begin{gate}[No hidden level-two term in this mother reduction]
For the particular ansatz of Ref.~\cite{AldrovandiSchaposnik2007APE4SQM},
the Chern--Simons density has no orientation-dependent term linear in
\(\dot n\); all time-dependent orientational terms in
Eq.~\eqref{eq:published-effective-action} come from the scalar kinetic
energy.  Thus this mother model derives \(C\) but does not derive a magnetic
line of Chern number two.  Other Chern--Simons vortex theories can induce
magnetic fields on moduli space~\cite{CollieTong2008APE4SQM}; that possibility
does not retroactively insert one here.
\end{gate}

\section{Physical tangent fermion and intrinsic \texorpdfstring{\(N=2\)}{N=2} SQM}

\subsection{BPS supersymmetry and tangent covariance}

The bulk vortex preserves one complex supercharge.  Supersymmetry therefore
completes Eq.~\eqref{eq:LB} with a complex Grassmann collective coordinate
\(\psi^w(t)\):
\begin{equation}
 L_{\rm SQM}=g_{w\bar w}\dot w\dot{\bar w}
 +\frac{\ii}{4}g_{w\bar w}
 \left(\bar\psi^{\bar w}D_t\psi^w-
 \bar D_t\bar\psi^{\bar w}\psi^w\right),
 \label{eq:sqm-action}
\end{equation}
with
\begin{align}
 D_t\psi^w&=\dot\psi^w+\Gamma^w{}_{ww}\dot w\psi^w,
 &\Gamma^w{}_{ww}&=g^{w\bar w}\partial_wg_{w\bar w}
 =-\frac{2\bar w}{1+|w|^2}.
 \label{eq:fermion-connection}
\end{align}
Under the south-chart coordinate \(v=1/w\),
\begin{equation}
 \psi^v=\frac{\partial v}{\partial w}\psi^w=-w^{-2}\psi^w,
 \qquad
 D_t\psi^v=-w^{-2}D_t\psi^w.
 \label{eq:fermion-transition}
\end{equation}
Hence \(\psi\in\Gamma(T^{1,0}\CP^1)\) and its kinetic norm is chart
independent.  This is the physical tangent fermion missing from a purely
bosonic AP-E4 argument.

\subsection{What the tangent fermion quantizes to}

Fix the canonical Spin\({}^c\) polarization associated with the complex
structure.  The tangent fermion and its adjoint become one Clifford
creation--annihilation pair.  Via the metric, a creation operator wedges by
a \((0,1)\)-form and its adjoint contracts.  Therefore the Fock bundle is
\begin{equation}
 W_{\rm can}=\Omega^{0,0}(\CP^1)\oplus\Omega^{0,1}(\CP^1)
 =\Omega^{0,*}(\CP^1),
 \qquad \Gamma=(-1)^q\quad\hbox{on }\Omega^{0,q}.
 \label{eq:fock-module}
\end{equation}
The supercharges and Hamiltonian are
\begin{equation}
 Q=\sqrt2\,\bar\partial,\qquad
 Q^\dagger=\sqrt2\,\bar\partial^\dagger,\qquad
 D=Q+Q^\dagger,\qquad
 H=\frac12\{Q,Q^\dagger\}=\frac12D^2.
 \label{eq:untwisted-supercharges}
\end{equation}
In particular \(Q^2=(Q^\dagger)^2=0\),
\(\{D,\Gamma\}=0\), and supersymmetric ground states are Dolbeault
cohomology classes.

The phrase ``tangent fermion'' describes the transformation of the
\emph{Clifford generator}.  It does not say that the Fock vacuum or every
wavefunction carries an additional factor \(T^{1,0}\CP^1\).  Declaring
\begin{equation}
 \Omega^{0,*}\longrightarrow
 \Omega^{0,*}\otimes T^{1,0}\CP^1
 \label{eq:forbidden-double-count}
\end{equation}
is an independent twist, not a consequence of
Eq.~\eqref{eq:fermion-transition}.

\begin{remark}[Vacuum-line ordering]
The explicit operator ordering used in the original vortex quantization
packages the two Fock components with a half-form factor and reports no
normalizable zero mode~\cite{AldrovandiSchaposnik2007APE4SQM}; geometrically
this is the ordinary-spin realization
\(\Omega^{0,*}\otimes K^{1/2}\), with
\(K^{1/2}=\cO(-1)\) on \(S^2\).  AP-E4's operator card instead declares the
canonical Spin\({}^c\) realization \eqref{eq:fock-module}.  The physical
tangent fermion is intrinsic, whereas the vacuum-line regularization is an
additional datum that must remain fixed when comparing spectra.  In
particular, the mother action does not derive the canonical vacuum line used
in the subsequent spectral theorem.  Neither choice supplies the desired
coefficient \(\cO(2)\).
\end{remark}

\section{Untwisted canonical negative controls and the independent
\texorpdfstring{\(\cO(2)\)}{O(2)} twist}

\subsection{The untwisted model does not have three vacua}

For the canonical realization, Hodge theory gives
\begin{equation}
 \Ker D^+\simeq H^0(\CP^1,\cO)=\CC,
 \qquad
 \Ker D^-\simeq H^1(\CP^1,\cO)=0.
 \label{eq:untwisted-count}
\end{equation}
Thus the declared canonical Spin\({}^c\) re-quantization has one ground
state.  As stressed in the vacuum-line remark, this is not the
source-selected half-form ordering, which has no normalizable zero mode.
A different model
with enough fermions to realize the full de Rham complex has
\(b_0+b_2=2\) vacua on \(S^2\).  This also is not a three-state mechanism.

\subsection{Magnetic/Wess--Zumino coefficient line}

Let \(E=\cO(k)\) be an independently derived Hermitian line.  In north and
south patches choose
\begin{align}
 a_N^{(k)}&=\frac{k}{2}(1-\cos\theta)\dd\phi,
 &a_S^{(k)}&=-\frac{k}{2}(1+\cos\theta)\dd\phi,
 \label{eq:monopole-connections}\\
 F^{(k)}&=\dd a^{(k)}=\frac{k}{2}\sin\theta\,
 \dd\theta\wedge\dd\phi,
 &\frac1{2\pi}\int_{S^2}F^{(k)}&=k.
 \label{eq:monopole-flux}
\end{align}
The overlap transition has winding \(k\), so integrality is necessary for a
global quantum line.  A worldline term
\begin{equation}
 S_{\rm WZ}=\hbar\int a^{(k)}_i(q)\dot q^i\dd t
 \label{eq:worldline-wz}
\end{equation}
and its supersymmetric curvature coupling replace
Eq.~\eqref{eq:untwisted-supercharges} by
\begin{equation}
 \cH_E=L^2\Omega^{0,*}(\CP^1,E),\qquad
 Q_E=\sqrt2\,\bar\partial_E,\qquad
 D_E=\sqrt2(\bar\partial_E+\bar\partial_E^\dagger),
 \qquad H_E=\frac12D_E^2.
 \label{eq:twisted-supercharges}
\end{equation}

For all integer \(k\), Riemann--Roch and Serre duality give
\begin{align}
 h^0(\cO(k))&=\max(k+1,0),
 &h^1(\cO(k))&=\max(-k-1,0),
 \label{eq:line-cohomology}\\
 \ind D_E^+&=h^0-h^1=k+1.
 \label{eq:rr-index}
\end{align}
At positive level two,
\begin{equation}
 \boxed{\dim\Ker D_{\cO(2)}^+=3,\qquad
 \dim\Ker D_{\cO(2)}^-=0,\qquad \ind D_{\cO(2)}^+=3.}
 \label{eq:three-states}
\end{equation}
A homogeneous-coordinate basis is
\(\{z_0^2,z_0z_1,z_1^2\}\).  The canonical Spin\({}^c\) determinant after
twisting is
\begin{equation}
 \det W_E=K^{-1}\otimes E^2=\cO(2k+2),
 \qquad k=2:\quad\det W_E=\cO(6).
 \label{eq:det-line}
\end{equation}
Orientation reversal is shifted in canonical Spin\({}^c\) convention:
three negative-chirality modes require \(E=\cO(-4)\), for which
\((h^0,h^1)=(0,3)\), not \(E=\cO(-2)\).

\begin{gate}[CPT and antiparticle polarization]
AP-E3 assigns \(k=-2\) to the \(B=-1\) antibaryon.  Holding the same
canonical complex polarization fixed would give
\((h^0(\cO(-2)),h^1(\cO(-2)))=(0,1)\), not the CPT conjugate of the
three-state \(B=+1\) kernel.  A physical completion must therefore derive how
CPT conjugates the complex/Spin\({}^c\) vacuum line---equivalently, why the
antiparticle operator is the anti-canonical conjugate, or why its effective
fixed-polarization coefficient is \(\cO(-4)\).  This polarization map is not
fixed by the signed WZW integer alone and remains an independent
composability gate.
\end{gate}

\section{Exact partner spectrum and spectral gap}

For \(k\ge0\), the monopole-harmonic multiplet at excitation number \(n\)
has
\begin{equation}
 j=\frac{k}{2}+n,\qquad \dim V_j=k+2n+1.
 \label{eq:j}
\end{equation}
The connection Laplacian and the K\"ahler--Weitzenb\"ock curvature shift
combine to
\begin{align}
 D_E^2
 &=\frac1{R^2}\left[j(j+1)-\frac{k^2}{4}-\frac{k}{2}\right]
 =\frac{n(n+k+1)}{R^2}
 =\frac{4n(n+k+1)}{C}.
 \label{eq:spectrum-square}
\end{align}
At \(n=0\) this reproduces \(k+1\) positive-chirality zero modes.  For
\(n\ge1\), oddness of \(D_E\) pairs the two chiral sectors and gives the
full spectrum~\cite{BykovSmilga2023APE4SQM}
\begin{equation}
 \boxed{\lambda_{n,\pm}=\pm\frac{2}{\sqrt C}
 \sqrt{n(n+k+1)},\qquad
 \operatorname{mult}(\lambda_{n,\pm})=k+2n+1.}
 \label{eq:full-spectrum}
\end{equation}
For \(k=2\),
\begin{equation}
 \Delta_D=\frac4{\sqrt C},\qquad
 \Delta_H=\frac12\Delta_D^2=\frac8C,
 \qquad \operatorname{mult}(\pm\Delta_D)=5.
 \label{eq:gaps}
\end{equation}
Every massive state has a partner.  Only the kernel carries a chiral
imbalance; therefore ``no partner'' must never be applied to the full tower.

\section{AP-E3/AP-E4 composability: necessary and sufficient data}

AP-E3 has produced an anomaly-consistent intermediate mixed-WZW candidate
whose signed collective term is intended to have \(k=n_cB=+2\).  That datum
can supply the line \(E=\cO(2)\) in Eq.~\eqref{eq:twisted-supercharges}, but
the numerical equality of two integers is not enough.

Let \(\Sigma_3\) be a spatial slice carrying the unit soliton,
\(\widetilde\cC_{B=1}\) the corresponding framed field-configuration
component before quotienting gauge transformations, and
\begin{equation}
 \operatorname{ev}:\Sigma_3\times\widetilde\cC_{B=1}\longrightarrow X
 \label{eq:evaluation-map}
\end{equation}
the evaluation map into the mixed-WZW target.  For the AP-E3 character
\(\widehat\Omega_5\in\widehat H^5(X;\ZZ)\), fiber integration in
differential cohomology first transgresses the degree-five WZW character to
the degree-two differential line class
\begin{equation}
 \widehat{\mathcal L}_{\rm WZW}
 :=\int_{\Sigma_3}\operatorname{ev}^*\widehat\Omega_5
 \in\widehat H^2(\widetilde\cC_{B=1};\ZZ).
 \label{eq:wzw-transgression}
\end{equation}
This degree-lowering step is mandatory: a degree-five class cannot be
directly pulled back and equated to a first Chern class.  Put
\(\cC_{B=1}:=\widetilde\cC_{B=1}/\cG\).  Gauge consistency further requires
the transgressed differential line to be basic and descend:
\begin{equation}
 \widehat{\mathcal L}_{\rm WZW}
 =p^*\widehat{\mathcal L}_{\rm WZW}^{\rm quot},
 \qquad p:\widetilde\cC_{B=1}\to\cC_{B=1}.
 \label{eq:wzw-gauge-descent}
\end{equation}
Failure of this descent is precisely a global gauge-anomaly obstruction.

Let \(\cM_{B=1}\) be the collective-coordinate component of the
charged-two-colour soliton and let
\begin{equation}
 \iota:\cM_{B=1}\longrightarrow\cC_{B=1}
 \label{eq:collective-embedding}
\end{equation}
embed it into the UV configuration space modulo gauge transformations.
The AP-E3 and AP-E4 cards are composable only if all of the following are
proved:
\begin{enumerate}[label=(C\arabic*)]
\item \(\cM_{B=1}\simeq\CP^1\) with a finite \(L^2\) metric and a positive
  gap to every noncollective fluctuation;
\item the projector \(P=nn^\dagger\) is the same physical orientation in
  both constructions, rather than an abstractly isomorphic sphere;
\item \(B[\iota(n)]=+1\) is constant on this component and the orientation
  map has degree \(+1\);
\item Eq.~\eqref{eq:wzw-gauge-descent} holds, including every large gauge
  transformation, so the transgressed line exists on the same quotient that
  contains the collective moduli;
\item the transgressed differential WZW line pulls back as
\begin{equation}
 \iota^*\widehat{\mathcal L}_{\rm WZW}^{\rm quot}
 =2\,\widehat c_1(\cO(1)),
 \qquad
 \frac1{2\pi}\int_{\cM_{B=1}}\iota^*F_{\rm WZW}=+2;
 \label{eq:signed-pullback}
\end{equation}
\item the pullback admits the supersymmetric completion in
  Eq.~\eqref{eq:twisted-supercharges} and introduces no extra massless Fermi or
  gauge modes;
\item CPT maps the \(B=+1\) canonical polarization to the correct
  anti-canonical \(B=-1\) operator, so the antibaryon has the conjugate
  three-state kernel rather than the one-state fixed-polarization
  \(\cO(-2)\) kernel;
\item the collective states carry the required microscopic representation,
  survive all gauge/anomaly projections, and couple through a Route-E portal
  below both the bulk and SQM gaps.
\end{enumerate}
When (C1)--(C8) hold, the composition is a theorem:
\begin{equation}
 k_{\rm eff}=n_cB\deg(\iota)=2\cdot1\cdot1=+2,
 \qquad E_{\rm eff}=\cO(+2),
 \label{eq:conditional-composition}
\end{equation}
and Eq.~\eqref{eq:three-states} follows.  If the degree is \(-1\), or if the
WZW line lives on a different sphere, the conclusion does not follow.

\begin{gate}[Mother-model mismatch]
The mother model in Sections 2--4 is a \(2+1\)-dimensional supersymmetric
\(U(2)\) Chern--Simons--Higgs vortex.  It independently proves that a
physical tangent fermion can arise on \(\CP^1\), but it is not the current
four-dimensional charged-two-colour
\(SU(2)_c\times U(1)_g\) AP-E3 candidate and does not prove
Eq.~\eqref{eq:collective-embedding}.  Therefore both the charged-two-colour
embedding flag and the AP-E4 worldline-completion flag remain false.
\end{gate}

Three microscopic ways to close the coefficient-line gate should be pursued
in parallel:
\begin{enumerate}
\item derive Eq.~\eqref{eq:signed-pullback} directly from the AP-E3
  mixed-WZW differential character on an explicit stable \(B=1\) solution;
\item calculate a bulk Callias/Fermi-zero-mode index over the soliton and
  show that its determinant line over \(\cM_{B=1}\) is \(\cO(+2)\);
\item derive an exactly-two gapped microscopic doublet whose Berry
  determinant over the \emph{same} moduli sphere has Chern number \(+2\),
  and prove that integrating it out preserves supersymmetry and the gap.
\end{enumerate}

\section{Backup product compactification and its six-dimensional anomaly gate}

A product construction
\(M^{3,1}\times\CP^1\) could instead turn the internal operator into a
four-dimensional chiral multiplicity via
\begin{equation}
 \slashed D_6=\slashed D_4\otimes1+\gamma_5\otimes D_E,
 \qquad \Gamma_7=\gamma_5\otimes\Gamma_2.
 \label{eq:product-dirac}
\end{equation}
This route is not used in the present proof.  It is admissible only after a
six-dimensional chiral field content passes all of the following:
\begin{enumerate}[label=(6D\arabic*)]
\item the complete local eight-form anomaly polynomial, including
  \(\Tr F^4\), \((\Tr F^2)^2\), \(p_1(T)\Tr F^2\), \(p_1(T)^2\), and
  \(p_2(T)\), vanishes or factorizes as
  \(I_8=\tfrac12\Omega_{\alpha\beta}X_4^\alpha X_4^\beta\) with an integral
  Green--Schwarz lattice~\cite{AlvarezGaumeWitten1984APE4SQM};
\item global gauge/gravitational anomalies are checked by the appropriate
  spin/Spin\({}^c\) bordism invariant rather than inferred from \(I_8\);
\item the Spin\({}^c\) determinant and \(\cO(2)\) flux obey integral and
  mod-two quantization conditions;
\item the index-weighted four-dimensional spectrum cancels every local and
  global anomaly, and the KK tower remains above the portal scale.
\end{enumerate}
None is established here, so
\texttt{six\_dimensional\_anomaly\_gate\_closed=false}.

\section{Deterministic verification}

The companion script
\path{route_f/code/verify_ap_e4_moduli_space_sqm.py} uses a fixed seed and
contains no fitted spectral data.  It verifies:
\begin{enumerate}
\item horizontal projection, the two-chart Fubini--Study metric, tangent
  fermion covariance, and invariant fermion norm at 250 deterministic
  complex probes;
\item positivity and convergence of the eliminated-\(\rho\) analytic
  radial surrogate;
\item Riemann--Roch and Serre duality for \(\cO(k)\), \(-6\le k\le6\),
  including the untwisted, \(k=2\), \(k=-4\), and de Rham controls;
\item midpoint curvature quadrature and transition-function winding for
  \(c_1(\cO(2))=2\);
\item a finite exact superalgebra block testing \(Q^2=0\), Hermiticity,
  \(\{D,\Gamma\}=0\), \(D^2=2H\), three positive-chirality zero modes,
  the first gap, and five states per massive sign;
\item fail-closed machine-readable fields separating intrinsic tangent
  physics from the charged-two-colour embedding, WZW pullback, Route-E
  portal, and six-dimensional anomaly gate.
\end{enumerate}
It writes
\begin{center}
\small
\path{route_f/output/ap_e4_moduli_space_sqm.json}\\
\path{route_f/output/ap_e4_moduli_space_sqm.md}.
\end{center}
A green run certifies the declared geometry and algebra; it deliberately does
not turn an open embedding into a physical theorem.

\section{Claim ledger and ordered next steps}

\begin{longtable}{@{}L{0.47\textwidth}L{0.18\textwidth}L{0.27\textwidth}@{}}
\toprule
Claim & Status & Boundary\\
\midrule
\endfirsthead
\toprule
Claim & Status & Boundary\\
\midrule
\endhead
Independent BPS mother has internal \(\CP^1\) and finite \(L^2\) metric &
\status{derived} & published radial reduction plus chart proof\\
Physical tangent fermion in intrinsic SQM & \status{derived} & BPS
collective-coordinate supersymmetry and Eq.~\eqref{eq:fermion-transition}\\
Canonical Spin\({}^c\)/Dolbeault algebra and spectrum & \status{solved} &
declared canonical vacuum-line polarization\\
Mother model selects that canonical vacuum line & \status{false} &
source-selected half-form ordering has no normalizable zero mode\\
Untwisted tangent SQM gives three vacua & \status{false} & one vacuum only in
the declared canonical re-quantization; source ordering has none and the de
Rham control has two\\
Independent \(\cO(+2)\) coefficient gives three positive zero modes &
\status{proved conditionally} & Riemann--Roch, Hodge theory, and spectrum\\
This mother model itself derives \(\cO(+2)\) & \status{false} & no
orientation-dependent first-order term in its reduction\\
Same charged-two-colour \(B=1\) \(\CP^1\) moduli map & \status{open} & no
stable AP-E3 soliton/embedding theorem\\
Signed transgressed AP-E3 WZW line is \(\cO(+2)\) on that same moduli &
\status{open} & Eq.~\eqref{eq:signed-pullback} not computed\\
CPT maps the \(B=-1\) sector to the conjugate three-state operator &
\status{open} & fixed canonical \(E=\cO(-2)\) has only one negative mode\\
Six-dimensional product anomaly gate & \status{open backup} & product route
unused; \(I_8\), GS, and bordism audits open\\
Three states are Route-E chiral families & \status{open} & representation,
portal, anomaly, and scale gates unresolved\\
\bottomrule
\end{longtable}

The ordered research programme is:
\begin{enumerate}
\item establish or rule out a stable charged-two-colour \(B=1\) soliton with
  \(\CP^1\) internal moduli and compute its full fluctuation gap;
\item transgress the signed mixed-WZW differential character along the
  spatial three-cycle, pull the resulting differential line to that exact
  family, and derive its supersymmetric worldline completion;
\item independently compute the soliton Fermi/Callias index and determinant
  line, providing a non-WZW cross-check of \(\cO(+2)\);
\item derive the vacuum-line regularization and its CPT/anti-canonical map
  rather than imposing the canonical polarization only in the baryon sector;
\item only after those agree, construct the Route-E portal below
  \(\min(\Delta_{\rm bulk},4/\sqrt C)\) and redo the anomaly ledger;
\item retain product compactification only if its complete 6D local and
  global anomaly gate can be closed.
\end{enumerate}

The final fail-closed status is
\begin{equation}
 \boxed{\begin{gathered}
 \texttt{physical\_tangent\_fermion\_in\_intrinsic\_sqm\_derived=true},\\
 \texttt{mother\_model\_vacuum\_line\_matches\_canonical\_spinc=false},\\
 \texttt{o2\_twist\_three\_positive\_zero\_modes=true},\\
 \texttt{cpt\_antibaryon\_polarization\_closed=false},\\
 \texttt{charged\_two\_colour\_soliton\_sqm\_embedding\_closed=false},\\
 \texttt{ap\_e4\_worldline\_completion\_closed=false},\\
 \texttt{route\_e\_portal\_closed=false},\\
 \texttt{six\_dimensional\_anomaly\_gate\_closed=false},\\
 \texttt{physics\_promotion\_allowed=false}.
 \end{gathered}}
 \label{eq:final-status}
\end{equation}

\bibliographystyle{unsrt}
\bibliography{route_f/tex/ap_e4_moduli_space_sqm}

\end{document}
